\documentclass[12pt,a4paper]{article}

\input{/home/tabaka/Documents/GitHub/Defs/defs.tex}

\setcounter{zadaniaRozdzial}{4}

\begin{document}
	
	\begin{center}
		\LARGE Karta 4
	\end{center}
	\vspace{1.5cm}
	
	\ZbiorZadanie{Wyznacz wszystkie wartości parametru $m$, dla których równanie
		$$(m+1)x^2-3mx+m+1=0$$
	ma dwa różne pierwiastki takie, że ich suma jest nie większa niż $2{,}5$.}
	
	\ZbiorZadanie{Rozwiąż równanie
		$$2\sin^2x-7\cos x-5=0$$
	w przedziale $\langle 0,2\pi\rangle$.}
	
	\ZbiorZadanie{Rozwiąż równanie
		$$\sin^2 2x-4\sin^2x+1=0$$
	w przedziale $\langle 0,2\pi\rangle$.}
	
	\ZbiorZadanie{Dany jest trójkąt ostrokątny $ABC$, w którym $|AC|=5$ i $|AB|=8$. Pole tego trójkąta jest równe $10\sqrt{3}$. Oblicz promień okręgu opisanego na tym trójkącie.}
	
	\ZbiorZadanie{Podstawa $AB$ trójkąta równoramiennego $ABC$ ma długość $8$ oraz $|\angle BAC|=30^\circ$. Oblicz długość środkowej $AD$ tego trójkąta.}
	
	\ZbiorZadanie{W trójkącie równoramiennym wysokość opuszczona na podstawę jest równa $36$, a promień okręgu wpisanego w ten trójkąt jest równy $10$. Oblicz długości boków tego trójkąta i promień okręgu opisanego na tym trójkącie.}
	
	\newpage
	
	\setcounter{zadaniaIndex}{0}
	
	\ZbiorZadanie{$m\in(-1,\frac{4-2\sqrt{17}}{5})\cup(\frac{4+2\sqrt{17}}{5},5)$}
	
	\ZbiorZadanie{$\cos x = -\frac{1}{2}, x\in\{\frac{2}{3}\pi,\frac{4}{3}\pi\}$}
	
	\ZbiorZadanie{$\sin^4x=\frac{1}{4}, x\in\{\frac{1}{4}\pi,\frac{3}{4}\pi,\frac{5}{4}\pi,\frac{7}{4}\pi\}$}
	
	\ZbiorZadanie{$R=\frac{7\sqrt{3}}{3}$. Wskazówka: Pole trójkąta rozpiętego na dwóch bokach i kącie między nimi, tw. cosinusów, tw. sinusów.}
	
	\ZbiorZadanie{$s=\frac{4\sqrt{21}}{3}$, trójkąty szczególne, tw. cosinusów.}
	
	\ZbiorZadanie{Boki: $30,39,39$. $R=39$. twierdzenie Pitagorasa, twierdzenie o odcinach stycznych/trójkąty podobne, tw sinusów.}
	
\end{document}